% Irodov, supplied PDF, printed p.50, PDF p.59.
% Last four problems of section 1.6, elastic deformations.
\exproblem{311}{50}
把一条长 $\ell=2.0\,\mathrm m$、
宽 $h=6.0\,\mathrm{cm}$、
厚 $\delta=2.0\,\mathrm{mm}$ 的钢带弯成闭合圆环。
在变形始终处于弹性范围、忽略端点连接损耗时，
至少需要做多少功？
\exsolution
成环半径为 $R=\ell/(2\pi)$。
设钢带宽度方向不发生弯曲，弯曲中性层沿带长，
厚度方向距中性层 $z$ 的纤维具有
轴向应变 $\varepsilon(z)=z/R$。
每单位体积的应变能密度为
$u(z)=E\varepsilon(z)^2/2$。
钢带体积元为
$\mathrm dV=h\ell\,\mathrm dz$，
所以总储能
\[
U=\frac{Eh\ell}{2R^2}
\int_{-\delta/2}^{\delta/2}z^2\,\mathrm dz
=\frac{Eh\ell\delta^3}{24R^2}.
\]
由 $R=\ell/(2\pi)$，
\[
\boxed{W_{\min}=U=
\frac{\pi^2Eh\delta^3}{6\ell}}.
\]
取钢材杨氏模量 $E\simeq2.0\times10^{11}\,\mathrm{Pa}$，
代入 $\ell,h,\delta$ 得
$\boxed{W_{\min}\simeq79\,\mathrm J\approx0.08\,\mathrm{kJ}}$。
\excheck 本结果是\emph{弯曲弹性能}，
不是将钢带整体均匀轴向伸长的能量；
弯曲应变沿厚度从负到正线性分布。

\exproblem{312}{50}
一根长 $\ell=1.0\,\mathrm m$、
半径 $r=10\,\mathrm{mm}$ 的钢质实心圆杆，
一端固定、另一端绕轴扭转
$\varphi=60^\circ$。
求杆内储存的扭转弹性形变能。
\exsolution
实心圆截面的极惯性矩为 $J=\pi r^4/2$。
由线弹性扭转关系 $N=GJ\varphi/\ell$，
从初始零角慢慢转到 $\varphi$ 时外力所做功为
\[
U=\int_0^\varphi N(\phi)\,\mathrm d\phi
=\frac{GJ\varphi^2}{2\ell}
=\boxed{\frac{\pi G r^4\varphi^2}{4\ell}}.
\]
采用钢的剪切模量
$G\simeq8.0\times10^{10}\,\mathrm{Pa}$、
$\varphi=\pi/3$、$r=10^{-2}\,\mathrm m$，
\[
\boxed{U\simeq6.9\times10^2\,\mathrm J}.
\]
\excheck \textbf{源书数值争议：}本版 PDF 题面明确是
$r=10\,\mathrm{mm}$、$\ell=1\,\mathrm m$，
书末虽印出同一正确公式，却给 $7\,\mathrm J$。
按所印半径与常用钢材 $G$ 实算为约 $690\,\mathrm J$，
相差约 $10^2$ 倍；若半径原拟为
$10/\sqrt{10}\,\mathrm{mm}$ 才近似 $7\,\mathrm J$，
但不能擅改原题数据。须将其作为待校数值错误记录。

\exproblem{313}{50}
一根长 $\ell$ 的均匀钢制实心圆杆，
一端固定，另一端绕杆轴扭转角 $\varphi$。
求杆内单位体积弹性形变能密度 $u(s)$
随离轴距离 $s$ 的分布。
\exsolution
在线弹性小扭转近似中，沿杆长角位移线性变化，
单位长度扭转角为 $\varphi/\ell$。
距转轴半径 $s$ 处的剪切应变
\[
\gamma(s)=\frac{s\varphi}{\ell}.
\]
由剪切胡克定律 $\tau=G\gamma$，
体积弹性能密度为从零应变加载的积分
\[
u(s)=\int_0^{\gamma(s)}G\gamma'\,\mathrm d\gamma'
=\frac12G\gamma(s)^2
=\boxed{\frac{G\varphi^2s^2}{2\ell^2}}.
\]
\excheck 轴心 $s=0$ 处没有剪切形变能，
外缘最大，且按半径的平方增长。
对全横截面及全长积分应恢复第312题公式。

\exproblem{314}{50}
求水深 $h=1000\,\mathrm m$ 处，
因液体受压所储存的\emph{单位体积弹性势能}，
水的密度为 $\rho$、等温体积压缩系数为 $\beta$。
\exsolution
以水面相对压强为零，忽略密度随深度的微小变化，
深度 $h$ 处静水相对压强为 $p=\rho gh$。
压缩系数定义为
$\beta=-(1/V)\mathrm dV/\mathrm dp$，
在小变形与常数 $\beta$ 近似下，
单位原始体积的压缩能为
\[
u=\int_0^{p}p'\beta\,\mathrm dp'
=\boxed{\frac12\beta(\rho gh)^2}.
\]
取水 $\beta\simeq4.8\times10^{-10}\,\mathrm{Pa^{-1}}$、
$\rho\simeq1000\,\mathrm{kg/m^3}$，得到
\[
\boxed{u\simeq2.3\times10^4\,\mathrm{J/m^3}
\approx23.5\,\mathrm{kJ/m^3}}.
\]
\excheck $\beta p^2$ 的量纲为
$\mathrm{Pa}=\mathrm{J/m^3}$；这是相对压强导致的弹性能，
没有把水的重力势能密度重复计入。
